%%%%%%%%%%%%%%%%%%%%%%%%%%%%%%%%%%%%%%%%%%%%%%%%%%%%%%%%%%%%%%%%%%%%%%
% Lớp 11 | Quan hệ vuông góc trong không gian | Hai đường thẳng vuông góc
% Dạng 1H8?1-1: Câu hỏi lý thuyết
% Số câu: 13  (sinh tự động bởi tach_cau.py; từ nay thầy sửa trực tiếp file này)
%%%%%%%%%%%%%%%%%%%%%%%%%%%%%%%%%%%%%%%%%%%%%%%%%%%%%%%%%%%%%%%%%%%%%%

%%%%%%%%%%%%%%%% TRẮC NGHIỆM 4 PHƯƠNG ÁN %%%%%%%%%%%%%%%%

% nguồn: Ok_Toan11_TracNghiem_Lan2
\begin{ex}%[1H8N1-1]%[ma:6cfad4a3]
	Trong không gian, cho đường thẳng $d$ và điểm $O$. Qua $O$ có bao  nhiêu đường thẳng vuông góc với đường thẳng $d$?
	\choice 
		 {$3$}
		{\True  vô số}
		 {$1$}
		 {$2$}
	\loigiai{ 
		Trong không gian, có  vô số đường thẳng qua một điểm cho trước và vuông góc với một đường thẳng cho trước.
	} 
\end{ex}

% nguồn: Ok_TracNghiemLop11_Lan1
\begin{ex}%[1H8N1-1]%[ma:3bafaff4]
	Chọn khẳng định đúng trong các  khẳng định sau:
	\choice 
		{\True Trong không gian, hai đường thẳng vuông góc với nhau có thể cắt nhau hoặc  chéo nhau}
		{Trong không gian, hai đường thẳng vuông góc với nhau thì phải cắt  nhau}
		{Trong không gian, hai đường thẳng không có điểm chung thì  song song với nhau}
		{Trong không gian, hai đường thẳng phân biệt cùng vuông góc với một đường  thẳng thì song song với nhau}
	\loigiai{ 
	 \immini{
		Giả sử cho hình  lập phương $ABCD.A_1B_1C_1D_1$.
		\begin{itemchoice}
			\itemch Đúng, ví  dụ $(AB,A_1D_1)=(AB, AD)=90^\circ$.
			\itemch Sai, ví dụ $AA_1\perp  DC$ nhưng $AA_1$ và $DC$ ở vị trí chéo nhau.
			\itemch Sai do  hai đường thẳng không có điểm chung thì song song hoặc chéo nhau.
			\itemch Sai, ví dụ $AD$ và $CC'$ cùng  vuông góc với $CD$ nhưng $AD$ không song song với $CC'$.
		\end{itemchoice}
		 }{
		\begin{tikzpicture}[scale=0.6,font=\footnotesize,line join=round,line cap=round,>=stealth]
				\coordinate  (a) at (2,2);
				\coordinate (b)  at (0,0);
				\coordinate  (c) at (5,0);
				\coordinate (d) at  (7,2);
				 \coordinate (a1) at ($(a)+(0,5)$);
				\coordinate (b1) at  ($(b)+(0,5)$);
				\coordinate (c1)  at ($(c)+(0,5)$);
				\coordinate (d1)  at ($(d)+(0,5)$);
				\draw[dashed] (b) -- (a)  -- (d)
					(a) -- (a1); 
				\draw (b) -- (c) -- (d) -- (d1) --  (c1) -- (b1) -- (a1) -- (d1)
					 (b) -- (b1)
					(c) -- (c1); 
				 \draw (a) node[below]{$A$};
				\draw  (b) node[below]{$B$};
				\draw (c) node[below]{$C$}; 
				\draw (d)  node[right]{$D$};
				\draw (a1)  node[above]{$A_1$};
				\draw (b1) node[below  left]{$B_1$};
				\draw (c1) node[below  right]{$C_1$};
				\draw  (d1) node[right]{$D_1$};
			\end{tikzpicture}
		} 
	 }
\end{ex}

% nguồn: Ok_TracNghiemLop11_Lan1
\begin{ex}%[1H8N1-1]%[ma:e25b25a8]
	Chọn khẳng định đúng trong  các khẳng định sau?
	 \choice
		 {Trong không gian hai đường thẳng phân biệt cùng vuông góc với một đường thẳng thì song song với nhau}
		{\True Trong  không gian hai đường thẳng vuông góc với nhau có thể cắt nhau hoặc chéo nhau}
		{Trong không gian hai mặt phẳng cùng vuông góc với một đường thẳng  thì song song với nhau}
		{Trong không gian hai đường thẳng không có  điểm chung thì song song với nhau}
	 \loigiai{
		"Trong không gian hai đường thẳng vuông góc với nhau có thể cắt nhau hoặc chéo nhau." là khẳng  định đúng.
	} 
\end{ex}

% nguồn: Ok_Toan11_TracNghiem_Lan2
\begin{ex}%[1H8H1-1]%[ma:ba4c8007]
	Chọn khẳng định đúng trong các  khẳng định sau:
	\choice 
		{Trong không  gian hai đường thẳng phân biệt cùng vuông góc với một đường thẳng thì song song với nhau}
		{\True  Trong không gian hai đường thẳng vuông góc với nhau có thể cắt nhau hoặc chéo nhau}
		{Trong không gian hai mặt phẳng cùng vuông góc với  một đường thẳng thì song song với nhau}
		{Trong  không gian hai đường thẳng không có điểm chung thì song song với nhau}
	 \loigiai{
		\begin{itemize}
			\item Trong không gian hai đường thẳng phân biệt cùng vuông góc với một đường thẳng có thể cắt nhau hoặc chéo nhau. \\\textbf{Ví dụ:} Cho hình lập phương $ABCD.A'B'C'D'$ ta có $\heva{&AA'\perp AB \\&AD\perp AB}$. Ta thấy $AA'$ và $AD$ cắt  nhau.
			\item Trong không  gian hai mặt phẳng cùng vuông góc với một đường thẳng có thể trùng nhau.
			\item Trong không  gian hai đường thẳng không có điểm chung thì có thể chéo nhau.
		\end{itemize}
	} 
\end{ex}

% nguồn: Ok_Toan11_TracNghiem_Lan2
\begin{ex}%[1H8H1-1]%[ma:102636cb]
	 Trong không gian, cho các mệnh đề sau, mệnh đề nào là mệnh đề đúng?
	\choice 
		{Một đường thẳng vuông góc với một trong hai đường thẳng vuông góc thì vuông góc với đường thẳng còn  lại}
		{Hai đường thẳng cùng song song với đường thẳng thứ ba thì song song với nhau} 
		{\True Một đường thẳng vuông góc với một trong hai đường thẳng song  song thì vuông góc với đường thẳng còn lại}
		{Hai đường  thẳng cùng vuông góc với đường thẳng thứ ba thì vuông góc với nhau}
	 \loigiai{
	 Theo định lí ta có $\heva{&a\perp b \\&b \parallel c}\Rightarrow a\perp c$.}
\end{ex}

% nguồn: Ok_TracNghiemLop11_Lan1
\begin{ex}%[1H8H1-1]%[ma:47f185f0]
	Trong không gian cho ba đường  thẳng phân biệt $a$, $b$, $c$. Khẳng định nào sau đây đúng?
	 \choice
		{Nếu đường thẳng $a$ vuông góc với đường thẳng $b$ và đường thẳng $b$ vuông góc với đường thẳng $c$  thì đường thẳng $a$ vuông góc với đường thẳng $c$}
		{\True Nếu đường thẳng $a$ vuông góc với đường thẳng $b$ và đường thẳng $b$ song song  với đường thẳng $c$ thì đường thẳng $a$ vuông góc với đường thẳng $c$}
		{Trong không gian, nếu đường  thẳng $a$ song song với đường thẳng $b$ và đường thẳng $b$ vuông góc với đường thẳng $c$ thì đường thẳng $a$ cắt đường thẳng $c$ tại một điểm}
		{Nếu ba đường thẳng $a$, $b$, $c$ vuông góc với nhau từng đôi một và có đường thẳng  $d$ vuông góc với đường thẳng $a$ thì đường thẳng $d$ song song với $b$ hoặc $c$}
	 \loigiai{
	Theo cách xác định góc giữa hai đường thẳng trong không gian, vì $b \parallel c$ nên ta có $(a,b)=(a,c)=90^\circ$ hay $a \perp c$. Suy ra chọn đáp án B.} 
\end{ex}

% nguồn: Ok_TracNghiemLop11_Lan1
\begin{ex}%[1H8H1-1]%[ma:c6b24012]
	Cho hai đường thẳng $a, b$ và mặt phẳng $(P)$. Mệnh đề  nào sau đây {\bf sai}?
	 \choice
		{Nếu $a\parallel (P)$ và $b\bot  (P)$ thì $a\bot b$}
		{Nếu $a\subset  (P)$ và $b\bot (P)$ thì $a\bot b$}
		{\True Nếu $a\parallel (P)$ và $b\bot a$ thì $b\bot (P)$} 
		{Nếu $a\bot (P)$ và $b\bot a$ thì $b\parallel (P)$ hoặc $b\subset (P)$} 
	\loigiai{ 
		Nếu $a\parallel (P)$ và  $b\bot a$ thì $b\bot (P)$ là một mệnh đề sai.
	 }
\end{ex}

% nguồn: Ok_TracNghiemLop11_Lan1
\begin{ex}%[1H8H1-1]%[ma:e74a969b]
	Mệnh đề  nào dưới đây là đúng?
	 \choice
		{Nếu một đường thẳng vuông góc  với một trong hai đường thẳng song song thì cũng vuông góc với đường thẳng còn lại}
		{Hai đường thẳng vuông góc với nhau thì luôn  cắt nhau}
		{Góc giữa hai đường thẳng bằng góc giữa hai  vec-tơ chỉ phương của chúng}
		{\True Hai đường thẳng cùng vuông góc với một đường thẳng thì song song với  nhau}
	 \loigiai{
		\begin{itemchoice}
			\itemch Đúng. Giả sử $a \perp b$ và $b \parallel c$ thì $(a,b)=(a,c)=90^\circ$, vậy  $a \perp c$.
			\itemch Sai vì hai đường thẳng vuông góc  có thể chéo nhau.
			\itemch Sai vì góc  giữa hai đường thẳng luôn nhỏ hơn hoặc bằng còn góc giữa hai vectơ có thể là góc tù.
			\itemch Sai vì hai đường thẳng cùng  vuông góc với một đường thẳng thì có thể cắt nhau.
		\end{itemchoice}
	 }
\end{ex}

% nguồn: Ok_TracNghiemLop11_Lan1
\begin{ex}%[1H8H1-1]%[ma:90ba7b1b]
	Mệnh đề nào sau đây {\bf  đúng}?
	\choice 
		{Nếu hai đường thẳng vuông góc với nhau thì hai đường thẳng đó cắt nhau} 
		{Hai đường thẳng cùng vuông góc với một đường  thẳng thì vuông góc với nhau}
		{Hai đường thẳng cùng vuông góc với một đường thẳng thì song song với nhau} 
		{\True Một đường  thẳng vuông góc với một trong hai đường thẳng song song thì vuông góc với đường thẳng kia}
	 \loigiai{
		Một đường  thẳng vuông góc với một trong hai đường thẳng song song thì vuông góc với đường thẳng kia.
	} 
\end{ex}

% nguồn: Ok_TracNghiemLop11_Lan1
\begin{ex}%[1H8H1-1]%[ma:bd2758c9]
	Trong  các mệnh đề sau, mệnh đề nào sai?
	 \choice
		{\True Hai đường thẳng phân biệt cùng vuông góc với một đường  thẳng thứ ba thì song song với nhau}
		{Hai đường thẳng phân biệt  cùng song song với một đường thẳng thì song song với nhau}
		{Một đường thẳng vuông góc với hai cạnh của tam giác thì sẽ vuông góc với cạnh  thứ ba của tam giác đó}
		{Hai đường thẳng vuông góc nếu góc giữa hai véc tơ chỉ  phương của chúng bằng $90^\circ$}
	 \loigiai{
		Hai đường thẳng phân biệt cùng vuông góc với đường thẳng thứ ba có thể song song, cắt  nhau hoặc chéo nhau.
	} 
\end{ex}

% nguồn: Ok_TracNghiemLop11_Lan1
\begin{ex}%[1H8B1-1]%[ma:2b4ad784]
	Chọn mệnh đề \textbf{đúng} trong các mệnh đề sau đây: 
	\choice 
		{Qua một điểm có duy nhất một mặt phẳng vuông góc với một mặt phẳng cho trước} 
		{\True Cho hai đường thẳng chéo nhau $a$ và $b$  đồng thời $a\perp b$. Luôn có mặt phẳng $(\alpha)$ chứa $a$ và $(\alpha) \perp b$}
		{Cho hai đường thẳng $a$ và $b$ vuông góc với nhau. Nếu mặt phẳng $(\alpha)$ chứa $a$ và mặt phẳng $(\beta)$ chứa $b$ thì  $(\alpha)\perp(\beta)$}
		{Qua một đường thẳng có duy nhất một  mặt phẳng vuông góc với một đường thẳng khác}
	 \loigiai{
		Có vô số mặt phẳng đi qua một điểm và vuông góc với một mặt phẳng cho trước. Do đó, A sai. Nếu hai đường thẳng $a$ và $b$ vuông góc với nhau và cắt nhau  thì mặt phẳng chứa cả $a$ và $b$ không thể vuông góc với $b$. Do đó, C sai.
		Qua một đường thẳng có vô số mặt phẳng  vuông góc với một đường thẳng khác. Do đó, D sai.
	 }
\end{ex}

% nguồn: Ok_TracNghiemLop11_Lan1
\begin{ex}%[1H8Y1-1]%[ma:3dcbdc6a]
	Trong các mệnh đề sau, mệnh đề nào  \textbf{sai}?
	 \choice
		{Hai đường thẳng cùng vuông góc với một mặt phẳng thì song  song}
		{Hai mặt phẳng phân biệt  cùng vuông góc với một đường thẳng thì song song}
		{\True Hai đường thẳng phân biệt cùng  vuông góc với một đường thẳng thì song song}
		{Hai mặt phẳng phân biệt cùng vuông góc với một mặt phẳng thì hai  mặt phẳng đó song song hoặc cắt nhau}
	 \loigiai{
		Trong không gian, hai đường thẳng phân biệt cùng vuông  góc với một đường thẳng thì có thể song song hoặc chéo nhau.\\
		Đáp án C chỉ đúng trong  mặt phẳng.
	 }
\end{ex}


%%%%%%%%%%%%%%%% ĐÚNG/SAI %%%%%%%%%%%%%%%%

% nguồn: Ok_dungsai_lop11_Lan1
\begin{ex}%[1H8N1-1]%[ma:49ae99c5]
 Trong không gian, cho ba đường  thẳng phân biệt. Xác định tính đúng sai các mệnh đề sau?
  \choiceTF
  {\True Nếu $a \parallel  b$ thì $(a,c)=(c,b)$}
  {\True Nếu $c \parallel b$ thì $(a,b)=(a,c)$} 
  {Nếu  $a \perp c$, $b \perp c$ thì $a \parallel b$}
   {Nếu $a \perp c$ thì $(a,b)=(c,b)$}
 \loigiai{ 
  \begin{itemchoice}
   \itemch Đúng vì theo định nghĩa góc giữa hai đường  thẳng trong không gian là góc giữa đường thẳng lần lượt song song hoặc trùng với chúng.
   \itemch Đúng vì theo định nghĩa góc giữa hai đường thẳng trong không gian là góc giữa đường thẳng lần lượt song song hoặc  trùng với chúng.
   \itemch Sai  vì $a$ và $b$ có thể chéo nhau.
   \itemch Sai vì theo  định nghĩa góc giữa hai đường thẳng trong không gian là góc giữa đường thẳng lần lượt song song hoặc trùng với chúng.
 \end{itemchoice}}
\end{ex}
